\section{Safety / AE / SAE}

% SLIDE ID: sec08-goal
\BlockGoalFrame{\SlideID{sec08-goal}}{
  \item zrozumieć przepływ informacji o bezpieczeństwie między pacjentem, ośrodkiem, \texttt{EDC} i systemem safety,
  \item odróżnić \texttt{AE}, \texttt{SAE} i \texttt{SUSAR},
  \item zobaczyć, gdzie najczęściej pojawiają się opóźnienia i błędy.
}

% SLIDE ID: sec08-ae-sae-susar
\begin{frame}{AE, SAE i SUSAR}
  \SlideID{sec08-ae-sae-susar}
  \begin{itemize}
    \item \texttt{AE}: każde niekorzystne zdarzenie medyczne u uczestnika badania.
    \item \texttt{SAE}: zdarzenie ciężkie, np. hospitalizacja, zagrożenie życia, zgon.
    \item \texttt{SUSAR}: ciężka, nieoczekiwana reakcja związana z badanym lekiem.
    \item Kluczowe pytanie brzmi nie tylko: \emph{co się stało?}, ale też: \emph{jak szybko i gdzie to trafi?}
  \end{itemize}
\end{frame}

% SLIDE ID: sec08-safety-flow
\begin{frame}{Przepływ safety od objawu do follow-up}
  \SlideID{sec08-safety-flow}
  \begin{enumerate}
    \item Pacjent zgłasza objaw.
    \item Ośrodek dokumentuje \texttt{AE}.
    \item Badacz ocenia ciężkość i związek z leczeniem.
    \item Pytanie decyzyjne: czy to \texttt{SAE}?
    \item Jeśli tak, rusza pilne raportowanie do sponsora / \texttt{CRO}.
    \item System safety prowadzi dalszy follow-up i medical review.
  \end{enumerate}
\end{frame}

% SLIDE ID: sec08-edc-dependency
\begin{frame}{Dlaczego safety zależy od danych w EDC}
  \SlideID{sec08-edc-dependency}
  \begin{itemize}
    \item \texttt{EDC} dostarcza kontekst: daty wizyt, leki, wyniki badań, kryteria włączenia.
    \item Dane z \texttt{EDC} pomagają ocenić czas wystąpienia zdarzenia i jego tło kliniczne.
    \item Niespójność między \texttt{AE} w \texttt{EDC} a zgłoszeniem safety tworzy ryzyko inspekcyjne.
    \item Jakość wpisu w \texttt{EDC} wpływa na szybkość i jakość medical review.
  \end{itemize}
\end{frame}

% SLIDE ID: sec08-medical-review
\begin{frame}{Medical Review: gdzie dane zamieniają się w decyzję}
  \SlideID{sec08-medical-review}
  \begin{itemize}
    \item Medical reviewer sprawdza kliniczny sens zgłoszenia i jego kompletność.
    \item Weryfikowane są: ciężkość, związek z leczeniem, oczekiwalność i potrzeba follow-up.
    \item Przy niejasnościach wracają pytania do ośrodka lub do zespołu data management.
    \item To miejsce, w którym słabe dane operacyjne szybko stają się problemem regulacyjnym.
  \end{itemize}
\end{frame}

% SLIDE ID: sec08-sae-reporting
\begin{frame}{SAE: raportowanie i terminy}
  \SlideID{sec08-sae-reporting}
  \begin{itemize}
    \item \texttt{SAE} wymaga pilnej eskalacji po uzyskaniu informacji przez ośrodek.
    \item System powinien wspierać szybkie zgłoszenie w oknie wymaganym przez procedurę badania.
    \item Każde opóźnienie trzeba umieć wyjaśnić i udokumentować.
    \item Proces nie kończy się na pierwszym zgłoszeniu: potrzebne są aktualizacje i follow-up.
  \end{itemize}
\end{frame}

% SLIDE ID: sec08-edc-integration
\begin{frame}{Integracja EDC z systemem safety}
  \SlideID{sec08-edc-integration}
  \begin{itemize}
    \item Integracja może przekazywać dane o pacjencie, zdarzeniu i tle klinicznym.
    \item Ogranicza podwójne wpisywanie tych samych informacji.
    \item Wspiera spójność między \texttt{EDC}, safety database i raportowaniem.
    \item Nadal potrzebna jest kontrola człowieka: integracja nie zastępuje oceny medycznej.
  \end{itemize}
\end{frame}

% SLIDE ID: sec08-roles
\begin{frame}{Rola badacza i rola sponsora / CRO}
  \SlideID{sec08-roles}
  \begin{columns}[T,onlytextwidth]
    \column{0.48\textwidth}
    \begin{DefinitionBox}{Badacz / ośrodek}
      \begin{itemize}
        \item rozpoznaje i dokumentuje zdarzenie,
        \item ocenia ciężkość i związek,
        \item zgłasza \texttt{SAE} w terminie,
        \item odpowiada na pytania follow-up.
      \end{itemize}
    \end{DefinitionBox}
    \column{0.48\textwidth}
    \begin{ExampleBox}{Sponsor / CRO}
      \begin{itemize}
        \item odbiera i weryfikuje zgłoszenie,
        \item prowadzi medical review,
        \item nadzoruje terminy i jakość procesu,
        \item ocenia obowiązki dalszego raportowania.
      \end{itemize}
    \end{ExampleBox}
  \end{columns}
\end{frame}

% SLIDE ID: sec08-typical-errors
\begin{frame}{Typowe błędy w procesie safety}
  \SlideID{sec08-typical-errors}
  \begin{itemize}
    \item wpisanie \texttt{AE} do \texttt{EDC}, ale bez równoległego zgłoszenia \texttt{SAE},
    \item niespójne daty między dokumentacją źródłową, \texttt{EDC} i safety database,
    \item brak oceny związku z leczeniem albo błędna klasyfikacja ciężkości,
    \item opóźnione follow-up lub brak domknięcia otwartych pytań.
  \end{itemize}
\end{frame}

% SLIDE ID: sec08-demo
\DemoFrame{Przepływ AE / SAE w systemach}{
  \begin{itemize}
    \item gdzie w \texttt{EDC} pojawia się informacja o zdarzeniu,
    \item jaki moment uruchamia zgłoszenie do systemu safety,
    \item jak wygląda różnica między pierwszym zgłoszeniem a follow-up,
    \item gdzie widać odpowiedzialność ośrodka, sponsora i zespołu safety.
  \end{itemize}
}

% SLIDE ID: sec08-case-hospitalization
\CaseStudyFrame{Pacjent trafia do szpitala po wizycie baseline}{
  Pacjent dzień po wizycie \texttt{Baseline} trafia do szpitala z nasilonymi objawami. Ośrodek dowiaduje się o hospitalizacji telefonicznie, a część danych pojawia się później w dokumentacji medycznej.
}{
  Co musi się wydarzyć w systemach? Jak dokumentujemy \texttt{AE}, kiedy rozpoznajemy \texttt{SAE}, jakie dane trafiają do \texttt{EDC}, co idzie do safety database i kto pilnuje follow-up?
}

% SLIDE ID: sec08-wrap
\begin{frame}{Najważniejsza myśl z bloku safety}
  \SlideID{sec08-wrap}
  \begin{itemize}
    \item Safety to nie pojedynczy formularz, tylko zsynchronizowany proces.
    \item Największe ryzyko powstaje na styku: pacjent, ośrodek, \texttt{EDC}, system safety.
    \item Dobra technologia pomaga, ale terminowość i odpowiedzialność ludzi pozostają kluczowe.
  \end{itemize}
\end{frame}
